\section{Notation and proof roadmap}\label{sec:notation-roadmap}

\subsection{Core notation}

The symbols used repeatedly in the proof are collected here.  The band
level $D(i)$ is included explicitly so that the later references to the
$D=0$ and $D=2$ strata are self-contained.

\begingroup
\renewcommand{\arraystretch}{1.16}
\noindent\begin{tabularx}{\linewidth}{@{}>{$}l<{$}>{\raggedright\arraybackslash}X@{}}
\toprule
\text{Symbol} & \text{Meaning}\\
\midrule
p,\ k,\ f & $p\geq5$ is prime; $f$ is an ordinary level-one form of weight
$4\leq k<p-1$.\\
\theta,\ \omega_{p^a} & The operator $q\,d/dq$ and weight filtration modulo
$p^a$, respectively.\\
i=np+i' & Band index, with $n$ the segment number and $i'$ the position
inside the segment.\\
P & The full Hasse drop $p(p-1)$.\\
D(i) & The least integer $D$ such that
$k+2i+DP\geq\omega_p(\theta^{i-p}f)+p(p+1)$; this is the inherited band
upper-bound level.\\
w_- & The candidate lower weight $k+2i-P$ in the $D=0$ stratum.\\
A,\ S,\ Z & $A=E_{p-1}$ is the chosen integral lift of the Hasse invariant,
$S=\operatorname{div}(A\bmod p)$, and $Z=pS$ is its length-$p$ thickening.\\
Q_r & The Hasse quotient
$H^0(Z,\omega^r|_Z)\simeq
M_r(\mathbb F_p)/A^pM_{r-P}(\mathbb F_p)$.\\
\mathcal B_i,\ \mathcal L_i & The Bockstein class of $\theta^if$ at weight
$w_-$ and its shifted leading-block candidate.  Their difference is the
boundary residual.\\
X & The divided Eisenstein series $E_2/12$ modulo $p$; powers of $X$ form
the direct block basis used in the reduction.\\
\alpha,\ \Phi_r,\ \eta_r & The first Hasse jet $(A-1)/p$, the periodicity
quotient $(\theta^{r+p-1}f-\theta^rf)/p$, and the lift error
$(G_r-\theta^rf)/p$, all reduced modulo $p$.\\
Z_d & The block direction $X^d\theta^{p-k}f$ in the factorial reduction.\\
i^\dagger & The complementary index $p^2+1-k-i$ used in the Tate-dual
comparison.\\
\bottomrule
\end{tabularx}
\endgroup

\paragraph{Terminology.}
A \emph{first Hasse jet} is a class on the thickened Hasse divisor $Z$.
A \emph{Bockstein class} is the intrinsic obstruction to an integral
weight-lowering lift.  A \emph{minimal-weight projection} is a chosen
characteristic-zero representative of the smallest candidate weight.  The
\emph{block basis} is the direct basis built from powers of $X$ and theta
iterates.  A \emph{boundary residual} is the difference between the
Bockstein class and its shifted leading-block candidate on one of the three
affine boundary lines.

\subsection{Dependency roadmap}

The proof has five logically separate layers.  The first three are uniform
in the admissible parameters and do not use the conjectural scalar identity.

\begingroup
\renewcommand{\arraystretch}{1.16}
\small
\noindent\begin{tabularx}{\linewidth}{@{}>{\raggedright\arraybackslash}p{0.15\linewidth}>{\raggedright\arraybackslash}p{0.31\linewidth}>{\raggedright\arraybackslash}X@{}}
\toprule
Layer & Main result & Logical role\\
\midrule
Geometric & Perfect Hasse-quotient duality and theta skew-adjointness &
Uses duality on $Z$; independent of all finite computations and block
formulas.\\
Integral & Bockstein realization, Hasse inclusion, and the explicit
periodic first-Hasse-jet formula & Uses integral Eisenstein lifts and the
first ordinary low-point congruence; supplies the minimal-weight projection.\\
Filtration & Uniform separation of all three boundary directions and
nonvanishing of the factorial scalar candidate & Uses exact mod-$p$
filtration; independent of the factorial block reduction.\\
Computational & Exact classification of $692$ cross-wedge cells and
$2{,}324$ third-line stress tests & Proved only on the deposited finite
range; discovers and falsifies formulas but is not extrapolated to all $p$.\\
Conditional & Factorial block reduction and two triangular endpoint
coefficients & Would identify the actual boundary residuals and imply the
shifted-conic filtration law.\\
\bottomrule
\end{tabularx}
\endgroup

Thus the dependency chain is
\[
\begin{aligned}
 \boxed{\text{geometry}}+\boxed{\text{integral first Hasse jet}}
 &\Longrightarrow \boxed{\text{explicit Bockstein class}},\\
 \boxed{\text{explicit Bockstein class}}
 &\xRightarrow[\text{open}]{\text{block reduction}}
 \boxed{\text{shifted-conic law}}.
\end{aligned}
\]
The boundary-separation theorem branches from the filtration layer before
the open implication and therefore remains unconditional.
